\section{Guidelines}
\label{sec:implication}

We summarize some guidelines to suggest better practices for leveraging LLMs in unit testing.

\textbf{(1) Potential of LLM-based approaches.} 
As discussed in Section~\ref{sec:rq1}, our study shows that fine-tuning LLMs significantly outperforms state-of-the-art methods such as $ATLAS$ for assertion generation, \textsc{AthenaTest} for test generation, and \textsc{Ceprot} for test evolution across nearly all metrics. This improvement highlights the potential of LLMs for a wide range of unit testing tasks. Therefore, LLM-based approaches should be regarded as a primary strategy for achieving superior results.

\textbf{(2) Incorporating Additional Post-Processing.}
As detailed in Section~\ref{sec:rq1-test_generation}, LLMs 
fall short compared to A3Test, which, though also based on fine-tuning, incorporates a post-processing step to ensure naming consistency, parentheses completion, and test signature accuracy. This suggests that even a straightforward post-processing step can effectively enhance LLMs' performance. Thus, exploring post-processing techniques, especially those that integrate program analysis knowledge, may be essential to enhance LLM effectiveness further.

\textbf{(3) Importance of Input Length.}
In Section~\ref{sec:rq1-test-evolution}, we describe reallocating input space in the \textsc{Ceprot} approach by removing the edit sequence to include additional content—such as the original method, updated method, and original test. This adjustment led to a 4.7\% improvement in EM and a 24.1\% increase in CodeBLEU, demonstrating that extending input context length significantly enhances LLMs' performance in the future.

\textbf{(4) Selection of LLMs.}
As discussed in Section~\ref{sec:rq2}, we provide empirical recommendations for selecting LLMs in unit testing tasks. First, large-scale LLMs consistently outperform smaller ones, indicating that larger models should be prioritized when sufficient computing resources are available. The CodeLlama and DeepSeek-Coder series are especially worth considering for their strong performance. Second, encoder-decoder architectures perform better among models with fewer than one billion parameters. 
Thus, in resource-constrained settings, encoder-decoder models, particularly the CodeT5p series, maybe the optimal choice for unit testing tasks.

\textbf{(5) Potential Issue with Failing Tests and Build Errors.}
As illustrated in Fig.~\ref{fig:rq2_3_all_tasks}, a considerable proportion of generated test cases fail when executed against production code, regardless of model size.
Moreover, according to Table~\ref{tab:tg_comparisons_llm}, even the top-performing model, DeepSeek-Coder-6b, still suffers from a high failure and build error rate of 30.36\% on Defects4J.
These failures, often resulting from incorrect test prefixes or assertions, significantly hinder the effectiveness of automated test case generation and impair the precision of bug detection. 
Therefore, reducing failing and build error test rates is crucial for further advancements.


\textbf{(6) Runtime-based Metrics rather than Text-based Metrics.}
As discussed in Section~\ref{sec:metrics_comparison}, LLMs may perform poorly on text-based metrics due to the diversity of test cases, which significantly impacts test generation tasks. The other two tasks may also be affected. Therefore, it is recommended to use runtime-based metrics to evaluate approaches for unit testing tasks.


